\documentclass[a4paper]{article}
\usepackage[utf8]{inputenc}
\usepackage[spanish]{babel}
\usepackage{datetime}

% https://tex.stackexchange.com/a/212264/273585
\newdateformat{monthyeardate}{\monthname[\THEMONTH] de \THEYEAR}

\begin{document}

\thispagestyle{empty}  % prevent page number on first page
                        % https://latex.org/forum/viewtopic.php?t=25462

Francisco Figari - \textit{Curriculum Vitae} - \monthyeardate\today

\begin{itemize}

  \item Antecedentes docentes
  \begin{itemize}

    \item
      Ayudante de segunda en la materia Métodos Numéricos @ \textit{DC, FCEyN,
      UBA}.  \newline
      4 cuatrimestres ejercidos entre 2016 y 2018 distribuidos entre la
      práctica y el laboratorio.

    \item
      Clases particulares de matemática (secundario y CBC, entre 3 y 5 clases
      por alumno, alrededor de 5 alumnos).

  \end{itemize}

  \item Antecedentes científicos
  \begin{itemize}

    \item
      Pasante de investigación @ \textit{DC, FCEyN, UBA} entre agosto 2020 y
      julio 2021 bajo el programa de becas para pasantías de iniciación en la
      investigación. \newline
      Estudio sobre el uso eye tracking en navegadores web aplicado a análisis
      clínicos a distancia, realizado en torno a tesis de licenciatura.

    \item
      Pasante de investigación @ \textit{ARSLab, Carleton University} entre
      febrero y junio de 2018. \newline
      Estudio sobre la aplicabilidad de simulación al desarollo de firmware de
      sistemas embebidos.

  \end{itemize}

  \item Experiencia laboral
  \begin{itemize}

    \item
      Programador fullstack @ \textit{DubbingDigital} entre marzo 2021 y
      diciembre 2021.  \newline
      Extensión y mantenimiento de interfaz y servidores web de software de
      gestión.

    \item
      Programador lead @ \textit{UrbanSim} entre julio 2018 y noviembre 2020.
      \newline
      Inicialmente extensión de generador procedural de vivienda asequible
      multifamiliar luego desarrollo de plataforma web envolvente.

    \item
      Programador freelance @ \textit{LIAA, FCEyN, UBA} durante primer mitad
      de 2018.  \newline
      Exploración del uso de eye tracking al estudio de atención en navegadores
      web.

    \item
      Programador junior @ \textit{Eryx} durante 4 meses a principios de 2016.

    \item
      Ayudante de cocina y mozo @ \textit{La Capelina, San Fernando} durante
      veranos 2012 y 2013.

  \end{itemize}

  \item Skills laborales
  \begin{itemize}

    \item
      Experiencia profesional con múltiples lenguajes (\texttt{Javascript},
      \texttt{C}\verb|++|, \texttt{Python}) y paradigmas (objetos, eventos,
      funcional). \newline
      Siempre dispuesto a aprender nuevas maneras de programar.

    \item
      Disposición y preferencia a contribuir en equipos interdisciplinarios.
      \newline
      Experiencia trabajando con programadores, gerentes, científiques,
      docentes, arquitectes y urbanistas. \newline
      Dirección de equipos de hasta 3 programadores.

    \item
      Habilidad para construir interfaces web y servicios de backend. \newline
      Comprensión intermedia de la API \texttt{JavaScript} para interactuar con
      el navegador. \newline
      Conocimientos elementales sobre cómo levantar, mantener y migrar bases de
      datos relacionales.

    \item
      Experiencia manipulando objetos geoespaciales, ya sea para dibujar con
      ellos en frontend, para almacenarlos adecuadamente o para operar con
      ellos dentro de otros algoritmos. \newline
      Uso de APIs de geocodificación.

    \item
      Interés en aplicar y estudiar buenas prácticas de programación y
      desarrollo de software. \newline
      Conocimiento vago pero diario de sistemas \texttt{Unix}.

    \item
      Inglés (uso diario a través de estudio y ocio) y Francés (escolaridad).

  \end{itemize}

  \item Estudios
  \begin{itemize}

    \item
      Licenciatura en Ciencias de la Computación @ \textit{DC, FCEyN, UBA}
      iniciada en 2013 y apuntando a finalizar durante 2022.

    \item
      Secundario @ \textit{Licée Jean Mermoz, Buenos Aires}.

  \end{itemize}

\end{itemize}

\end{document}
